% Source of paper.pdf, the PDF that the import workflow annotates:
%   pdflatex paper.tex
% Each page starts with a black bar: half the paper width on page 1, a
% quarter of it on page 2.
\documentclass[12pt]{article}
\usepackage[letterpaper]{geometry}
\usepackage{amsmath,amssymb,amsthm}
\newtheorem{theorem}{Theorem}
\pagestyle{empty}
\begin{document}
\noindent\rule{0.5\paperwidth}{1cm}

\section*{Bases}

\begin{theorem}
Every vector space $V$ over a field $k$ has a basis.
\end{theorem}

\begin{proof}
Order the linearly independent subsets of $V$ by inclusion. The union of a
chain of such subsets is linearly independent, so Zorn's lemma gives a
maximal linearly independent subset $B$. If $v \in V$ were outside the span
of $B$, then $B \cup \{v\}$ would be linearly independent, against the
maximality of $B$. Hence $B$ spans $V$.
\end{proof}

\newpage
\noindent\rule{0.25\paperwidth}{1cm}

\section*{Dimension}

\begin{theorem}
Any two bases of $V$ have the same cardinality.
\end{theorem}

\begin{proof}
For finite bases, exchange the vectors of one basis for the vectors of the
other, one at a time: each exchange keeps a spanning set, so neither basis
has more vectors than the other.
\end{proof}
\end{document}
